\documentclass[12pt]{article}
\usepackage{amsmath,amssymb}
\usepackage[margin=1in]{geometry}
\usepackage{hyperref}

\begin{document}

\section{Overview}

This document summarizes the refinements considered for Panel B. I divide the
results into relatively good and bad results. A good result does not
necessarily mean that the specification should be used as the main result. It
means that the coefficient moves in the expected negative direction without
an obvious data-quality concern. Throughout the search, I exclude restrictions
based on discrepancies between the scraped data and Inside Airbnb data.

The original Panel B estimate is approximately $0.09^{***}$. The clearest
general pattern is that conditioning on $Superhost_{t-2}=1$ often moves the
coefficient toward zero or into negative territory. Price trimming alone does
not generate this pattern.


\section{Good Results}

\subsection{Conditioning on $Superhost_{t-2}=1$}

The most interpretable refinement is to restrict Panel B to listings that held
Superhost status two periods earlier, i.e., \texttt{ex\_super2 = ``t''}. This
constructs a more homogeneous group of listings that recently lost Superhost
status. Across the nine price and price-change trimming combinations, six
coefficients are negative, although none is statistically significant.

\begin{table}[h]
\centering
\caption{Panel B with \texttt{ex\_super2 = ``t''}: selected trimming results}
\label{tab:exsuper2_trim}
\begin{tabular}{lcccc}
\hline
Price trim & Change trim & $N$ & Conventional & Bias-corrected/robust \\
\hline
0\% & 0\% & 1,327 & $-0.065$ $(0.091)$ & $-0.067$ $(0.104)$ \\
1\% & 0\% & 1,304 & $-0.084$ $(0.087)$ & $-0.087$ $(0.100)$ \\
1\% & 1\% & 1,269 & $-0.005$ $(0.070)$ & $-0.005$ $(0.078)$ \\
5\% & 0\% & 1,128 & $-0.090$ $(0.112)$ & $-0.114$ $(0.130)$ \\
5\% & 1\% & 1,101 & $-0.050$ $(0.084)$ & $-0.062$ $(0.096)$ \\
\hline
\end{tabular}
\end{table}

This pattern suggests that \texttt{ex\_super2} is more important than the
trimming rule itself. Its main limitation is the decline in sample size. For
example, under the original 5\% price-change trimming rule, Panel B falls from
18,483 observations to 1,181 observations after applying \texttt{ex\_super2}.


\subsection{Rating and review activity after applying \texttt{ex\_super2}}

After applying \texttt{ex\_super2}, a rating cutoff of 4.7 or above makes the
FULL, Q1Q2, Q2Q3, and Q3Q4 coefficients negative. However, none is
statistically significant, and the sample falls to 818 observations. The main
pattern is therefore favorable, but the cutoff should be justified before it
is used.

Review activity produces stronger negative estimates. The main candidates are
reported below. These estimates use 5\% price trimming and no price-change
trimming.

\begin{table}[h]
\centering
\caption{Review-based refinements of Panel B}
\label{tab:review_refinements}
\small
\begin{tabular}{p{0.38\textwidth}ccc}
\hline
Restriction & $N$ & Conventional & Bias-corrected/robust \\
\hline
Prior reviews $\geq$ quarter median
 & 515 & $-0.159^{**}$ $(0.076)$ & $-0.181^{**}$ $(0.085)$ \\
Prior reviews $\geq 30$
 & 561 & $-0.142^{**}$ $(0.071)$ & $-0.152^{*}$ $(0.080)$ \\
Prior reviews $\geq$ quarter median and recency $\leq$ quarter median
 & 248 & $-0.102$ $(0.065)$ & $-0.093$ $(0.073)$ \\
\hline
\end{tabular}
\end{table}

The quarter-median rule is less arbitrary than an absolute review cutoff and
produces a statistically significant negative coefficient. However, its
effect is larger than the preferred 0.15 bound. The joint median restriction
satisfies the effect-size bound, but it is not statistically significant.
The 30-review result is also a boundary case: the conventional estimate is
within the bound, while the bias-corrected estimate is slightly outside it.


\subsection{Conditioning on recent listing activity}

Conditioning on \texttt{ex\_quarter\_ltm $\geq 5$} without
\texttt{ex\_super2} changes the FULL coefficient to $-0.017$ $(0.026)$ and
produces several negative lower-price estimates. This is an improvement over
the original positive coefficient, but it remains statistically insignificant
and the cutoff is somewhat arbitrary. The higher-price subgroup also does not
follow the Panel A pattern.


\subsection{Panel A consistency check}

As a consistency check, I reconstructed the quarterly data using 5\% price
trimming and no price-change trimming. For Panel A, all six FULL-sample
specifications remain negative and statistically significant. The conventional
estimates are $-0.105^{***}$, $-0.045^{**}$, $-0.094^{***}$,
$-0.053^{***}$, $-0.039^{**}$, and $-0.208^{***}$. Thus, the revised data
construction does not remove the main Panel A pattern.


\section{Bad Results}

\subsection{Price and price-change trimming without \texttt{ex\_super2}}

Trimming the price level and the price change at 0\%, 1\%, or 5\% produces
nine combinations. None generates a negative Panel B estimate. The estimates
range from approximately zero to 0.108, and most are positive and statistically
significant. Therefore, trimming alone does not solve the Panel B problem.

\begin{table}[h]
\centering
\caption{Panel B estimates from trimming alone}
\label{tab:trim_only}
\begin{tabular}{lccc}
\hline
 & \multicolumn{3}{c}{Price-change trimming} \\
Price trimming & 0\% & 1\% & 5\% \\
\hline
0\% & $0.060^{*}$ & $0.000$ & $0.090^{***}$ \\
1\% & $0.078^{**}$ & $0.056^{**}$ & $0.108^{***}$ \\
5\% & $0.089^{***}$ & $0.063^{**}$ & $0.092^{***}$ \\
\hline
\end{tabular}
\end{table}


\subsection{Past running-variable restrictions}

I considered restrictions such as $x_{t-1}\geq 4.5$ and
$x_{t-1}\geq 4.75$. The first-stage estimates become unstable across Panel B
subgroups, while the second-stage coefficients remain positive. These
restrictions are therefore not useful.


\subsection{Covariate cutoffs with little improvement}

Conditioning on \texttt{ex\_quarter\_rating $\geq 4.5$}, cumulative reviews
$\geq 5$, or \texttt{ex\_quarter\_ltm $\geq 1$ or $\geq 3$} produces results
similar to the original specification. The Panel B coefficients remain
positive in the main sample. Combining LTM $\geq 3$ with cumulative reviews
$\geq 5$ also produces a positive FULL estimate of 0.038 $(0.025)$.

After applying \texttt{ex\_super2}, LTM cutoffs from 1 to 5 do not improve the
FULL estimate either: the estimates remain positive, between 0.024 and 0.079.
Combining LTM and rating cutoffs reduces the sample further and generally moves
the estimates toward zero rather than producing stable negative effects.


\subsection{Other listing filters}

The following filters do not produce a stable, statistically significant FULL
effect:
\begin{itemize}
    \item non-missing host response rate: 0.011 $(0.054)$, $N=891$;
    \item days since last review $\leq 60$: 0.002 $(0.069)$, $N=378$;
    \item minimum nights $\geq 30$ and cumulative reviews $\leq 10$:
    $-0.097$ $(0.106)$, $N=256$.
\end{itemize}
The last two restrictions also leave very small samples, and their absolute
cutoffs require additional justification.


\subsection{Large and unstable exploratory estimates}

Restricting the sample to single-listing hosts produces a conventional
estimate of $-0.254^{*}$ and a bias-corrected estimate of $-0.277^{*}$ under
one trimming rule. Combining the single-listing restriction with rating
$\geq 4.7$ produces even larger estimates, $-0.334^{*}$ and $-0.376^{*}$.
These magnitudes are well outside the 0.15 effect-size bound. They also do not
survive correction for the large number of specifications considered, so I do
not treat them as credible main results.


\section{Overall Pattern}

The results point to three conclusions. First, \texttt{ex\_super2} is the only
simple restriction that repeatedly moves Panel B estimates in the expected
negative direction. Second, price trimming by itself does not explain the
positive Panel B coefficient. Third, review-based filters can produce negative
and sometimes significant estimates, but the significant estimates either
slightly exceed the preferred effect-size bound or depend on a cutoff that is
difficult to justify. The most defensible bounded result is the joint
quarter-median review and recency restriction, but it is not statistically
significant.

These findings should therefore be described as exploratory. No clean
specification examined so far is both statistically significant under
conventional and robust inference and below 0.15 in absolute value. Results
based on cross-source data discrepancies are intentionally excluded.

\vspace{0.5em}
\noindent\footnotesize
\textit{Notes:} Unless otherwise stated, reported results correspond to the
first RD specification: side-specific MSE-optimal bandwidths, time fixed
effects, triangular kernel, fuzzy RD at 4.75, and standard errors clustered by
listing ID. $^{***}p<0.01$, $^{**}p<0.05$, $^{*}p<0.10$.

\end{document}
